\section{Geometric consequences}\label{sec:consequences}

Theorem~\ref{thm:main} has consequences for other notions of dimension,
for unions of segments with prescribed directions, and for several
curved incidence problems. The direction-set and curved conclusions
use the transfer theorems of Keleti and M\'ath\'e, Gao, Liu, and Xi,
and Nadjimzadah.

\subsection{Packing and Minkowski dimensions}

Every Kakeya set \(K\subset\R^4\) also has packing dimension four, since
Hausdorff dimension is at most packing dimension, which is at most
the ambient dimension. If the set is bounded, both its lower and
upper Minkowski dimensions are four: Hausdorff dimension is at most
lower Minkowski dimension, and lower Minkowski dimension is at most
upper Minkowski dimension, which is at most four. Here the lower and
upper Minkowski dimensions are respectively the lower and upper
limits of \(\log N_\delta(K)/\log(1/\delta)\), where
\(N_\delta(K)\) is the least number of balls of radius \(\delta\)
covering \(K\), as \(\delta\downarrow0\). Thus the theorem also
controls the common-scale covering growth of each bounded Kakeya set.

\subsection{Projection to four dimensions}

Theorem~\ref{thm:main} also implies that every Kakeya set
\(E\subset\R^d\), \(d\ge4\), has \(\dim_H E\ge4\); in particular,
this holds in dimension five. Fix a four-dimensional linear subspace
\(V\), and let \(P_V\) be orthogonal projection onto \(V\). For each
unit \(v\in V\), a witnessing segment \(a+[0,1]v\subset E\) projects
to the unit segment \(P_Va+[0,1]v\). Hence \(P_VE\) is a Kakeya set
in \(V\simeq\R^4\), so \(4=\dim_H(P_VE)\le\dim_H E\). The final
inequality follows because projection does not increase covering
diameters, with no compactness or measurability assumption.
The resulting lower bound is four; the full conjecture in
\(\R^d\) asks for dimension \(d\).

\subsection{Arbitrary directions and extension of segments}

Keleti and M\'ath\'e transfer the full-direction assertion to arbitrary
direction sets. We identify unoriented directions in \(\R^4\) with
the real projective space \(\mathbb{RP}^3\).

\begin{corollary}[General direction-set bound]
\label{cor:general-directions-four}
Let \(D\subset\mathbb{RP}^3\) be nonempty and let \(E\subset\R^4\).
If \(E\) contains a nondegenerate line segment in every direction
in \(D\), then
\[
 \dim_H E\ge1+\dim_H D.
\]
\end{corollary}
\begin{proof}
By \citet[Theorem~1.4]{KeletiMatheEquivalences}, there is a compact
Besicovitch set \(B\subset\R^4\) such that
\(\dim_H E\ge\dim_H B-(3-\dim_H D)\).
Their Besicovitch convention requires a unit segment in every
direction, so Theorem~\ref{thm:main} gives \(\dim_H B=4\).
Substitution proves the bound.
\end{proof}

In particular, a direction set of Hausdorff dimension three already
forces full spatial dimension. We will use this below when the
directions form a nonempty open patch. Neither \(E\) nor \(D\) is
required to be measurable, and the positive segment lengths may vary.

The same conclusion gives a line-segment extension principle by the
implication proved in \citet[Section~1.2]{KeletiMatheEquivalences}.
For any nonempty family \(\mathcal S\) of nondegenerate segments in
\(\R^4\), let \(\mathcal L\) be the family of their supporting full
lines. Then
\[
 \dim_H\Bigl(\bigcup_{S\in\mathcal S}S\Bigr)
 =\dim_H\Bigl(\bigcup_{L\in\mathcal L}L\Bigr).
\]
Thus extending every segment in a fixed family to its whole line
preserves the Hausdorff dimension of the union. This application uses
the authors' same-dimensional implication from the general
direction-set bound.

\subsection{Nikodym sets on constant-curvature manifolds}

The Nikodym formulation selects geodesics through a positive-volume
family of base points, rather than selecting a segment for every
direction. Theorem~\ref{thm:main} controls this different incidence
requirement through the established Kakeya-to-Nikodym implication of
Gao, Liu, and Xi. We use
their normalized local geodesic formulation
\citep[Definition~1.3 and Theorem~2.1]{GaoLiuXi2025}: the metric is
normalized so that the injectivity radius is at least \(10\), and the
geodesic segments below have unit length. For a Borel set
\(\Omega\subset M\) and \(0<\lambda<1\), put
\[
 \Omega_\lambda^\star
 =\left\{x\in M:
 \begin{array}{l}
 \text{there is a unit geodesic segment \(\gamma_x\) through \(x\)}\\
 \text{such that }|\gamma_x\cap\Omega|_g\ge
          \lambda|\gamma_x|_g
 \end{array}\right\},
\]
where the bars denote geodesic length. The set \(\Omega\) is a
\emph{Nikodym set} if \(\Omega_\lambda^\star\) has positive
Riemannian volume for every \(\lambda\) sufficiently close to \(1\).

\begin{corollary}[Four-dimensional Nikodym dimension]
\label{cor:nikodym-four}
Let \((M^4,g)\) be a Riemannian manifold of constant sectional
curvature in the preceding normalized setting. Every Nikodym set
\(\Omega\subset M\) has \(\dim_H\Omega=4\).
\end{corollary}
\begin{proof}
Theorem~\ref{thm:main} applies in particular to the compact Euclidean
Kakeya sets in the convention of
\citet[Definition~1.1]{GaoLiuXi2025}. It therefore supplies their
Kakeya Hausdorff lower bound with \(d=\alpha=4\).
\citet[Theorem~1.5]{GaoLiuXi2025} transfers this bound to Nikodym sets
on constant-sectional-curvature \(d\)-manifolds, giving
\(\dim_H\Omega\ge4\). The reverse inequality holds for every subset
of a four-dimensional Riemannian manifold.
\end{proof}

This resolves Conjecture~1.4 of \citet{GaoLiuXi2025} in dimension
four and includes Euclidean \(\R^4\). The geodesic straightening and
the Kakeya-to-Nikodym transfer are results of those authors; no new
transfer argument or variable-curvature assertion is made here.

\subsection{Curved Kakeya sets for a fixed translation-invariant phase}
\label{subsec:curved-kakeya-four}

A separate consequence concerns the structured curved-Kakeya class of
Gao, Liu, and Xi. Unlike the preceding Nikodym corollary, it uses a
direction-indexed family from one phase in a Euclidean chart. We use
their phase and core conventions
\citep[Definitions~1.7--1.11]{GaoLiuXi2025}, with the active support
made explicit.
One common local change of coordinates straightens a direction patch
of this family. This lets the Euclidean dimension theorem control the
original curved set through local bi-Lipschitz invariance.

Write \(X=(x,t)\in\R^3\times\R\), and fix a sufficiently small chart
\[
 U=B_\rho^3\times(-\rho,\rho),\qquad Y=B_\rho^3,
\]
where \(B_\rho^3\) is the open Euclidean ball about the origin. Fix one
smooth phase on a neighborhood of \(\overline U\times\overline Y\),
\[
 \phi(x,t;y)=x\cdot y+\psi(t;y),\qquad \psi(0;y)=0.
\]
For \((X;y)\) in this neighborhood, let
\[
 G_0(X;y)=\bigwedge_{j=1}^3\partial_{y_j}\nabla_X\phi(X;y),
 \qquad \mathcal D_y=G_0(X;y)\cdot\nabla_X,
\]
where the wedge is identified with its normal vector in \(\R^4\),
and \(\mathcal D_y\) differentiates in \(X\) at fixed \(y\).
Impose the mixed-rank and curvature conditions
\[
 \begin{aligned}
 \rank\nabla_X\nabla_y\phi(X;y)&=3,\\
 \det\left(\left.\nabla_y^2
   \langle\nabla_X\phi(X;y),G_0(X;y_0)\rangle
   \right|_{y=y_0}\right)&\ne0,
 \end{aligned}
\]
and Bourgain's condition in the invariant formulation of
\citet[Theorem~2.1]{GuoWangZhangDichotomy},
\[
 \mathcal D_y^2\partial_{y_i y_j}^2\phi(X;y)
 =C(X;y)\mathcal D_y\partial_{y_i y_j}^2\phi(X;y)
 \quad(1\le i,j\le3),
\]
with the same scalar \(C(X;y)\) for every pair \(i,j\).
All three conditions hold throughout the chosen neighborhood of
\(\overline U\times\overline Y\), after one fixed phase-dependent
localization. Thus the active curve chart lies within the region of
validity of the support hypotheses; no condition is extended to
off-chart curve portions. For \(y\in Y\) and \(z\in U\), define the
full local core
\[
 \Gamma_y^\phi(z)
 =\{X\in U:\nabla_y\phi(X;y)=\nabla_y\phi(z;y)\}.
\]

\begin{corollary}[Fixed-chart curved Kakeya dimension]
\label{cor:curved-kakeya-four}
Under the preceding fixed-chart phase hypotheses, let
\(E\subset\R^4\). Suppose that for every \(y\in Y\) there exists
\(\omega_y\in B_\rho^3\) such that
\[
 \Gamma_y^\phi((\omega_y,0))\subset E.
\]
Then \(\dim_H E=4\).
\end{corollary}
\begin{proof}
For this phase, \(G_0=(-\partial_t\nabla_y\psi,1)\) in the standard
orientation. The mixed rank is automatic, and the remaining
conditions become
\[
 \det\nabla_y^2\partial_t\psi(t;y)\ne0,\qquad
 \nabla_y^2\partial_t^2\psi(t;y)
   =c(t,y)\nabla_y^2\partial_t\psi(t;y).
\]
The nonsingular matrix determines this scalar smoothly and
independently of \(x\).
\citet[Corollary~4.2 and Theorem~1.12]{GaoLiuXi2025} show that \(c\)
depends only on \(t\) on the connected local product and straighten
the phase. More precisely, their proof gives, after shrinking to an
open time interval \(I\) containing \(0\) and a nonempty direction
ball \(Y_0\subset Y\),
\[
 \kappa(u,s)=(u+B(\alpha(s)),\alpha(s)),\qquad
 \alpha(0)=0,\quad \alpha'(s)\ne0,
\]
\[
 \phi(\kappa(u,s);y)=u\cdot y+s h(y)+q(y)+f(s),\qquad
 \det\nabla_y^2h(y)\ne0,
\]
whenever \(s\in I\), \(y\in Y_0\), and \(\kappa(u,s)\in U\).
This is one smooth coordinate change for the fixed phase.

Since \(\psi(0;y)=0\), the core through \((\omega_y,0)\) is
\[
 \Gamma_y^\phi((\omega_y,0))
 =\{(\omega_y-\nabla_y\psi(t;y),t)\in U\}.
\]
Put \(V=U\cap(\R^3\times\alpha(I))\) and
\(F=\kappa^{-1}(E\cap V)\). Every anchor \((\omega_y,0)\) lies in the
open set \(V\), so continuity supplies a nontrivial closed time
interval on which its core remains in \(V\). In the new coordinates
the gradient equality defining that core is exactly
\[
 u=(\omega_y-B(0))-s\nabla h(y).
\]
Indeed, the \(\nabla q(y)\) terms cancel from the two gradient values,
and \(f(s)\) has zero \(y\)-gradient. Thus \(F\) contains a
positive-length straight segment in the direction
\[
 \nu(y)=\frac{(-\nabla h(y),1)}
               {\sqrt{1+|\nabla h(y)|^2}}
 \quad\text{for every }y\in Y_0.
\]
The nonsingular Hessian of \(h\) makes \(\nu\) a local diffeomorphism.
Its image therefore contains a nonempty open patch in the upper
hemisphere of \(S^3\). The corresponding unoriented directions form
an open set \(D\subset\mathbb{RP}^3\), so \(\dim_H D=3\).
Corollary~\ref{cor:general-directions-four} gives \(\dim_H F=4\);
it allows the positive segment lengths to depend on \(y\).

A smooth local diffeomorphism preserves Hausdorff dimension of
arbitrary subsets, using a countable cover by coordinate
neighborhoods on which it and its inverse are Lipschitz. Consequently
\[
 4=\dim_H F=\dim_H(E\cap V)\le\dim_H E\le4.
\]
No compactness or measurability of \(E\) or of \(y\mapsto\omega_y\),
and no uniform lower bound on the initial segment lengths, was used.
\end{proof}

The straightening and general-dimensional set transfer are due to
\citet[Theorem~1.12 and Section~4.2]{GaoLiuXi2025}. The corollary applies
their transfer to full gradient-defined cores in one fixed active
chart.

The translation-invariance assumption specifies the class being
straightened here. Bourgain's condition alone need not allow one coordinate change to
straighten the curves, as the examples of
\citet[Example~1.24 and Proposition~1.25]{NadjimzadahBourgain} show.

\subsection{A three-dimensional quadratic curved consequence}

A further application uses Nadjimzadah's lifting theorem to pass
from the four-dimensional linear result to curved sets in dimension
three. For a real \(2\times2\) matrix \(B\), let \(I_2\) denote the
identity matrix. Use the fixed compact anchor ball
\(W_B\subset\R^2\), chosen sufficiently large, and the fixed time
interval \(I_B=[-\tau_B,\tau_B]\), with \(\tau_B>0\) sufficiently
small, specified after equation~(1.16) of
\citet{NadjimzadahLifting}.

\begin{corollary}[Rank-one quadratic curved Kakeya sets]
\label{cor:quadratic-curved-three}
Let \(B\) be a real \(2\times2\) matrix with
\(\rank B=\rank(B^2)=1\), and use the fixed domains just specified.
Suppose that a compact set \(K\subset\R^3\) contains, for every
\(y\in[-1,1]^2\), the full curve
\[
 \{(w_y+(tI_2+t^2B)y,t):t\in I_B\}
\]
for some \(w_y\in W_B\). Then \(\dim_H K=3\).
\end{corollary}
\begin{proof}
Corollary~\ref{cor:general-directions-four} supplies the linear Kakeya
hypothesis in the convention of
\citet[Section~1.2, equations~(1.8)--(1.13)]{NadjimzadahLifting}.
Indeed, a compact linear-chart set \(F\subset\R^4\) in that convention
contains a segment
\[
 \{(b_\eta+t\eta,t):t\in J\}
\]
for every slope \(\eta\) in a fixed three-dimensional parameter ball,
with anchors in a fixed compact ball and one fixed nondegenerate
compact interval \(J\). Slopes in the interior give a nonempty open
set of directions \([\eta:1]\in\mathbb{RP}^3\), and every segment has
positive length. The corollary therefore gives \(\dim_H F=4\).

The geometry of Nadjimzadah's quadratic lift explains the two rank
hypotheses. Write \(B=UV\), where \(U:\R\to\R^2\) is injective and
\(V:\R^2\to\R\) is surjective, and define
\[
 \pi:\R^2\times\R\times\R\longrightarrow\R^2\times\R,
 \qquad \pi(x,z,t)=(x+tUz,t).
\]
The lifted segments in
\citet[equations~(4.1)--(4.9)]{NadjimzadahLifting} are
\[
 L_{y,w,r}(t)=\bigl((w,r)+t(y-Ur,Vy),t\bigr),
 \qquad t\in I_B,
\]
with \(r\) in a fixed compact interval with interior. They map to the
prescribed curves over the entire time interval:
\[
 \pi(L_{y,w,r}(t))=(w+ty+t^2UVy,t)
                  =(w+(tI_2+t^2B)y,t).
\]
Since \(B^2=U(VU)V\) has rank one, \(VU\ne0\). The slope map
\(\Theta(y,r)=(y-Ur,Vy)\) consequently satisfies
\[
 \det D\Theta
   =\det\begin{pmatrix}I_2&-U\\ V&0\end{pmatrix}=VU\ne0.
\]
Thus varying \(y\) and \(r\) in their interiors produces an open slope
patch in dimension four. Definition~2.1 of
\citet{NadjimzadahLifting} retains every \(t\in I_B\) in these lifted
segments, each of which has length at least \(2\tau_B\).

With the four-dimensional linear hypothesis established,
\citet[Theorem~1.3]{NadjimzadahLifting} applies with
\(3+\rank(B^2)=4\). It gives infimal Hausdorff dimension
\(3-\rank B+\rank(B^2)=3\) for the compact curved sets in the
statement. Hence \(\dim_H K\ge3\), and the ambient bound gives equality.
\end{proof}

This is a lower-dimensional application of the headline result:
the prescribed quadratic curves lift to linear Kakeya sets in one
additional dimension. Compactness and containment of the full curves
over the fixed interval are part of this application. No regularity
of the anchor choice \(y\mapsto w_y\) is required.
